\section{Detection, Thresholding, and Offloading}\label{sec:methods}

All detectors receive a colour image resized to $256\times256$ pixels, are fitted on $k$ defect-free images only, and return an anomaly map $A(\mathbf{x})$ that is bilinearly up-sampled from the feature grid and smoothed with a Gaussian ($\sigma=4$\,px). The image score is $s=\max_{\mathbf{x}}A(\mathbf{x})$.

\subsection{Station Model: Gabor--Gaussian}\label{sec:gg}
The image is converted to CIE L*a*b*. The lightness channel $L$ is filtered with quadrature Gabor pairs at wavelengths $\lambda\in\{4,8,16\}$\,px and orientations $\theta\in\{0^\circ,45^\circ,90^\circ,135^\circ\}$ (envelope $\sigma=0.56\lambda$, aspect ratio 0.5), giving 12 energy maps
\begin{equation}
E_{\lambda,\theta}=\sqrt{\left(L\star g^{\mathrm{e}}_{\lambda,\theta}\right)^2+\left(L\star g^{\mathrm{o}}_{\lambda,\theta}\right)^2}.
\end{equation}
All 24 correlations are computed with one FFT of the reflect-padded image and 12 inverse FFTs of precomputed complex spectra, which is numerically equivalent to spatial filtering (max.\ relative deviation $<10^{-6}$) but about $3\times$ cheaper on one core. For the 12 energies and for $L$, $a^*$, $b^*$, local mean and standard deviation over a $15\times15$ window are sampled with stride 4, yielding a $D=30$-dimensional descriptor $\mathbf{f}(\mathbf{x})$ on a $64\times64$ grid. Because fabric texture is approximately stationary, a \emph{single} Gaussian $\mathcal{N}(\boldsymbol{\mu},\boldsymbol{\Sigma})$ is fitted to the descriptors of all positions of all shots (at most $2\cdot10^5$ randomly chosen vectors), with ridge regularisation $\boldsymbol{\Sigma}_\epsilon=\boldsymbol{\Sigma}+\epsilon\,(\operatorname{tr}\boldsymbol{\Sigma}/D)\,\mathbf{I}$, $\epsilon=0.01$, and
\begin{equation}
A(\mathbf{x})=\sqrt{(\mathbf{f}(\mathbf{x})-\boldsymbol{\mu})^{\top}\boldsymbol{\Sigma}_\epsilon^{-1}(\mathbf{f}(\mathbf{x})-\boldsymbol{\mu})}.
\end{equation}
Sharing one Gaussian across positions means that even a single image supplies 4096 training vectors, so the model is usable at $k=1$. The profile consists of $D+D^2=930$ numbers plus the threshold.

\subsection{Edge Model: PatchCore-lite}\label{sec:pc}
Following PatchCore~\cite{roth2022patchcore}, ResNet-18 layer-2 ($128\times32\times32$) and layer-3 ($256\times16\times16$) features are averaged over $3\times3$ neighbourhoods, layer~3 is up-sampled to $32\times32$, and the two are concatenated into 384-dimensional patch descriptors. All descriptors of the $k$ shots form a memory bank that is randomly subsampled to 10\,\% (at least 1000 entries); a test patch is scored by its Euclidean distance to the nearest bank entry. Random rather than greedy-coreset subsampling keeps fitting essentially instantaneous; we therefore call the method PatchCore-lite. As a third, position-specific reference we also evaluate PaDiM~\cite{defard2021padim} (ResNet-18 layers 1--3, 100 random channels, one Gaussian per position), which is not deployed because its per-position statistics make the profile large and LOO calibration expensive.

\subsection{Label-Free Thresholds by Leave-One-Out}\label{sec:loo}
When a new fabric is learned, no defect examples exist. For each shot $i$ the detector is refitted on the other $k-1$ shots and the score $s_i^{\mathrm{LOO}}$ of shot $i$ is computed; these scores estimate the score distribution of \emph{unseen} defect-free frames (scoring the training shots with the full model would be biased low---for a nearest-neighbour model, trivially zero). The threshold is
\begin{equation}
\tau=(1+m)\max_{i=1..k}s_i^{\mathrm{LOO}},\label{eq:tau}
\end{equation}
with margin $m=0.1$ by default. LOO costs $k$ refits on cached features: milliseconds for the Gabor--Gaussian model and well under a second for PatchCore-lite at $k=10$. We denote the station and edge thresholds by $\tau_g$ and $\tau_p$.

\subsection{Confidence-Gated Offloading}\label{sec:gate}
Let $s_g$ be the station score of a frame. With $0<\alpha\le1\le\beta$, the station
\begin{equation}
\text{decision}=\begin{cases}
\text{OK (local)}, & s_g<\alpha\,\tau_g,\\
\text{DEFECT (local)}, & s_g>\beta\,\tau_g,\\
\text{offload}\;\Rightarrow\;\text{DEFECT iff } s_p>\tau_p, & \text{otherwise},
\end{cases}\label{eq:gate}
\end{equation}
where $s_p$ is the edge score. $\alpha=\beta=1$ reduces to station-only inspection; $\alpha\to0,\beta\to\infty$ to edge-only. Both thresholds come from (\ref{eq:tau}), so the band is expressed in units of the fabric's own normal-score scale and transfers across fabrics without labelled tuning. We fix $\alpha=0.6,\beta=1.5$ a priori and report the full $(\alpha,\beta)$ trade-off in Section~\ref{sec:res-gate}.

\subsection{Load Model}\label{sec:load}
Let $N$ stations capture $f$ frames/s, a fraction $o$ of which are offloaded with mean payload $J$ bytes (plus MQTT/TCP/IP overhead $h$). The per-station uplink rate and the edge utilisation are
\begin{equation}
R_{\mathrm{up}}=8fo(J+h),\qquad \rho=\frac{Nfo}{\mu(b)},\label{eq:load}
\end{equation}
where $\mu(b)$ is the edge throughput (frames/s) at batch size $b$. The edge can serve at most $N_{\max}=\mu/(fo)$ stations; cutting $o$ by gating raises $N_{\max}$ proportionally. The end-to-end latency of an offloaded frame is approximately
\begin{equation}
T\approx t_g+\frac{8(J+h)}{B}+2d+W_q+t_p(b),\label{eq:latency}
\end{equation}
with station compute $t_g$, uplink bit-rate $B$, one-way delay $d$, queueing delay $W_q$, and edge compute $t_p$; locally decided frames take $t_g$ only. Section~\ref{sec:res-live} compares (\ref{eq:load}) with measurements.
